% GENERATED FILE. Edit canonical lessons, then run npm run generate:books.
% canonical-source-hash: fnv1a64:b7abf5aa72e5d9eb
% canonical-lessons: AR-C394-atqana, AR-C394-iktasaba, AR-C394-taammaqa, AR-C394-lakhkhasa, AR-C394-mulakhkhas

\chapter{Learning well}
\label{ch:ar-learning-well}
\hlchaptermodality{394}

\begin{chapteropening}
\textbf{By the end of this chapter:} \emph{I can say that someone mastered, acquired, went deeper or summarized, and talk about a summary.}

Complete the last lesson of chapter 394: I can say that someone mastered, acquired, went deeper or summarized, and talk about a summary.
\end{chapteropening}

\section[\texorpdfstring{\ar{أتقن} (atqana)}{أتقن (atqana)}]{\ar{أتقن} (atqana) — to master}

\label{lesson:AR-C394-atqana}



\begin{quote}
Before the new one: say the Arabic for to excel, then the Arabic for to evaluate.
\end{quote}

\subsection*{You'll want to know: \ar{أتقن}}
\textbf{\ar{أتقن}} — \emph{atqana} — \textquotedblleft{}to master\textquotedblright{}.

\begin{quote}
\ar{أتقن} \ar{جدي} \ar{لغتين}.
\end{quote}

\begin{quote}
\emph{atqana jaddī lughatayn.}
\end{quote}

\begin{quote}
My grandfather mastered two languages.
\end{quote}

Present: \ar{يتقن}, \emph{yutqin}.

\begin{cousinweb}
\textbf{\ar{أتقن}} (\emph{atqana}) is built on the native Arabic root \emph{t-q-n}, \textquotedblleft{}to do a thing skilfully\textquotedblright{} (Lane). It is Form IV, to make a thing firm and exact: to master it. \ar{إتقان}, mastery, is its verbal noun.
\end{cousinweb}

\begin{grammarlens}[title={What you've built}]
One more word for giving reasons and opinions.
\end{grammarlens}

\subsection*{Your turn}
\begin{itemize}
\raggedright
  \item \emph{Say it:} \emph{atqana}
  \item \emph{Say it:} \emph{atqana}, once more
  \item \emph{Say it:} \emph{atqana jaddī lughatayn.} — My grandfather mastered two languages.
  \item \emph{Recall:} say the Arabic for to excel, then the Arabic for to evaluate, then say \emph{atqana} again
\end{itemize}

\begin{tcolorbox}[breakable,colback=teal!4,colframe=teal!35!black,title={Before you move on}]
What does \ar{أتقن} mean? (\textquotedblleft{}To master\textquotedblright{}.) Say it once more.
\end{tcolorbox}
\section[\texorpdfstring{\ar{اكتسب} (iktasaba)}{اكتسب (iktasaba)}]{\ar{اكتسب} (iktasaba) — to acquire}

\label{lesson:AR-C394-iktasaba}



\begin{quote}
Before the new one: say the Arabic for to evaluate, then the Arabic for to master.
\end{quote}

\subsection*{You'll want to know: \ar{اكتسب}}
\textbf{\ar{اكتسب}} — \emph{iktasaba} — \textquotedblleft{}to acquire\textquotedblright{}.

\begin{quote}
\ar{اكتسب} \ar{الموظف} \ar{خبرة} \ar{كبيرة}.
\end{quote}

\begin{quote}
\emph{iktasaba al-muwaẓẓaf khibra kabīra.}
\end{quote}

\begin{quote}
The employee gained a lot of experience.
\end{quote}

Present: \ar{يكتسب}, \emph{yaktasib}.

\begin{cousinweb}
\textbf{\ar{اكتسب}} (\emph{iktasaba}) is built on the native Arabic root \emph{k-s-b}, \textquotedblleft{}to gain, to earn\textquotedblright{} (Lane). It is Form VIII of \emph{kasaba}, to earn: to acquire for oneself.
\end{cousinweb}

\begin{grammarlens}[title={What you've built}]
One more word for giving reasons and opinions.
\end{grammarlens}

\subsection*{Your turn}
\begin{itemize}
\raggedright
  \item \emph{Say it:} \emph{iktasaba}
  \item \emph{Say it:} \emph{iktasaba}, once more
  \item \emph{Say it:} \emph{iktasaba al-muwaẓẓaf khibra kabīra.} — The employee gained a lot of experience.
  \item \emph{Recall:} say the Arabic for to evaluate, then the Arabic for to master, then say \emph{iktasaba} again
\end{itemize}

\begin{tcolorbox}[breakable,colback=teal!4,colframe=teal!35!black,title={Before you move on}]
What does \ar{اكتسب} mean? (\textquotedblleft{}To acquire\textquotedblright{}.) Say it once more.
\end{tcolorbox}
\section[\texorpdfstring{\ar{تعمق} (taʿammaqa)}{تعمق (taʿammaqa)}]{\ar{تعمق} (taʿammaqa) — to go deeper}

\label{lesson:AR-C394-taammaqa}



\begin{quote}
Before the new one: say the Arabic for to master, then the Arabic for to acquire.
\end{quote}

\subsection*{You'll want to know: \ar{تعمق}}
\textbf{\ar{تعمق}} — \emph{taʿammaqa} — \textquotedblleft{}to go deeper\textquotedblright{}.

\begin{quote}
\ar{تعمق} \ar{الطالب} \ar{في} \ar{دراسة} \ar{التاريخ}.
\end{quote}

\begin{quote}
\emph{taʿammaqa aṭ-ṭālib fī dirāsat at-tārīkh.}
\end{quote}

\begin{quote}
The student went deeper into the study of history.
\end{quote}

Present: \ar{يتعمق}, \emph{yataʿammaq}.

\begin{cousinweb}
\textbf{\ar{تعمق}} (\emph{taʿammaqa}) is built on the native Arabic root \emph{ʿ-m-q}, \textquotedblleft{}to be deep\textquotedblright{} (Lane). It is Form V, to go deep into a thing. \ar{عميق}, deep, is from the same root.
\end{cousinweb}

\begin{grammarlens}[title={What you've built}]
One more word for giving reasons and opinions.
\end{grammarlens}

\subsection*{Your turn}
\begin{itemize}
\raggedright
  \item \emph{Say it:} \emph{taʿammaqa}
  \item \emph{Say it:} \emph{taʿammaqa}, once more
  \item \emph{Say it:} \emph{taʿammaqa aṭ-ṭālib fī dirāsat at-tārīkh.} — The student went deeper into the study of history.
  \item \emph{Recall:} say the Arabic for to master, then the Arabic for to acquire, then say \emph{taʿammaqa} again
\end{itemize}

\begin{tcolorbox}[breakable,colback=teal!4,colframe=teal!35!black,title={Before you move on}]
What does \ar{تعمق} mean? (\textquotedblleft{}To go deeper\textquotedblright{}.) Say it once more.
\end{tcolorbox}
\section[\texorpdfstring{\ar{لخص} (lakhkhaṣa)}{لخص (lakhkhaṣa)}]{\ar{لخص} (lakhkhaṣa) — to summarize}

\label{lesson:AR-C394-lakhkhasa}



\begin{quote}
Before the new one: say the Arabic for to acquire, then the Arabic for to go deeper.
\end{quote}

\subsection*{You'll want to know: \ar{لخص}}
\textbf{\ar{لخص}} — \emph{lakhkhaṣa} — \textquotedblleft{}to summarize\textquotedblright{}.

\begin{quote}
\ar{لخص} \ar{المعلم} \ar{القصة} \ar{في} \ar{جملتين}.
\end{quote}

\begin{quote}
\emph{lakhkhaṣa al-muʿallim al-qiṣṣa fī jumlatayn.}
\end{quote}

\begin{quote}
The teacher summed up the story in two sentences.
\end{quote}

Present: \ar{يلخص}, \emph{yulakhkhiṣ}.

\begin{cousinweb}
\textbf{\ar{لخص}} (\emph{lakhkhaṣa}) is built on the native Arabic root \emph{l-kh-ṣ}, \textquotedblleft{}to make clear, to explain\textquotedblright{} (Lane). It is Form II; from making a thing clear it came to mean drawing out its gist: to summarize.
\end{cousinweb}

\begin{grammarlens}[title={What you've built}]
One more word for giving reasons and opinions.
\end{grammarlens}

\subsection*{Your turn}
\begin{itemize}
\raggedright
  \item \emph{Say it:} \emph{lakhkhaṣa}
  \item \emph{Say it:} \emph{lakhkhaṣa}, once more
  \item \emph{Say it:} \emph{lakhkhaṣa al-muʿallim al-qiṣṣa fī jumlatayn.} — The teacher summed up the story in two sentences.
  \item \emph{Recall:} say the Arabic for to acquire, then the Arabic for to go deeper, then say \emph{lakhkhaṣa} again
\end{itemize}

\begin{tcolorbox}[breakable,colback=teal!4,colframe=teal!35!black,title={Before you move on}]
What does \ar{لخص} mean? (\textquotedblleft{}To summarize\textquotedblright{}.) Say it once more.
\end{tcolorbox}
\section[\texorpdfstring{\ar{ملخص} (mulakhkhaṣ)}{ملخص (mulakhkhaṣ)}]{\ar{ملخص} (mulakhkhaṣ) — a summary}

\label{lesson:AR-C394-mulakhkhas}



\begin{quote}
Before the new one: say the Arabic for to go deeper, then the Arabic for to summarize.
\end{quote}

\subsection*{You'll want to know: \ar{ملخص}}
\textbf{\ar{ملخص}} — \emph{mulakhkhaṣ} — \textquotedblleft{}a summary\textquotedblright{}.

\begin{quote}
\ar{قرأت} \ar{ملخص} \ar{الكتاب}.
\end{quote}

\begin{quote}
\emph{qaraʾtu mulakhkhaṣ al-kitāb.}
\end{quote}

\begin{quote}
I read the summary of the book.
\end{quote}

\begin{cousinweb}
\textbf{\ar{ملخص}} (\emph{mulakhkhaṣ}) is built on the native Arabic root \emph{l-kh-ṣ}, \textquotedblleft{}to make clear, to explain\textquotedblright{} (Lane). \emph{Mulakhkhaṣ} is the Form II passive participle: what has been summed up, a summary.
\end{cousinweb}

\begin{grammarlens}[title={What you've built}]
One more word for giving reasons and opinions.
\end{grammarlens}

\subsection*{Your turn}
\begin{itemize}
\raggedright
  \item \emph{Say it:} \emph{mulakhkhaṣ}
  \item \emph{Say it:} \emph{mulakhkhaṣ}, once more
  \item \emph{Say it:} \emph{qaraʾtu mulakhkhaṣ al-kitāb.} — I read the summary of the book.
  \item \emph{Recall:} say the Arabic for to go deeper, then the Arabic for to summarize, then say \emph{mulakhkhaṣ} again
\end{itemize}

\begin{tcolorbox}[breakable,colback=teal!4,colframe=teal!35!black,title={Before you move on}]
What does \ar{ملخص} mean? (\textquotedblleft{}A summary\textquotedblright{}.) Say it once more.
\end{tcolorbox}
